\part{その他}
    \chapter{無限級数公式集}
        \section{\texorpdfstring{$\sum_{k=1}^{\infty}{\frac{\sin{k}}{k}} = \frac{\pi-1}{2}$}{Σ\_\{k=1\}\string^∞ sin(k)/k = (π-1)/2}}
            \begin{proof}
                \quad\par
                下図のような鋸波$y=f(x)$をフーリエ級数展開して求める。
                \begin{figure}[H]
                    \centering
                    \includegraphics[keepaspectratio, scale=0.5]
                    {parts/complexAnalysis/figs/series_frac{sin(n)}{n}/pic.pdf}
                    \caption{鋸波}
                \end{figure}
                $f(x)$は奇関数なので$\cos$成分はゼロ，つまり$\forall n, a_n=0$。$\sin(nx)$の係数は
                \[b_n=\frac{1}{\pi}\integrate{-\pi}{\pi}{x\sin(nx)}{}{x}=\frac{2}{\pi}\integrate{-\pi}{0}{x\sin(nx)}{}{x}=(-1)^{n-1}\frac{2}{n}\]
                よって
                \[x\sim2\left(\frac{\sin x}{1}-\frac{\sin 2x}{2}+\frac{\sin 3x}{3}-\dotsb\right)\]
                $x$を$x-\pi$で置き換えると
                \[x-\pi \sim -2\left(\frac{\sin x}{1}+\frac{\sin2x}{2}+\frac{\sin3x}{3}+\dotsb\right)\]
                特に$x$=1とすれば公式を得る。
            \end{proof}
        \section{\texorpdfstring{$\sum_{n=1}^\infty{\frac{1}{1+n^2}}=\frac{-1}{2}+\frac{\pi}{2}\coth{\pi}$}{Σ\_\{n=1\}\string^∞ 1/(1+n\string^2) = -1/2 + (π/2)coth(π)}}
            大阪大学 大学院 工学研究科 平成21年度 院試 数学 問題4 の解答を参照。
        \section{\texorpdfstring{$\sum_{n=1}^\infty{\frac{1}{(1+n^2)^2}}=\frac{-1}{2}+\frac{\pi^2}{4\sinh^2{\pi}}+\frac{\pi}{4}\coth\pi$}{Σ\_\{n=1\}\string^∞ 1/(1+n\string^2)\string^2 = -1/2 + π\string^2/(4sinh\string^2(π)) + (π/4)coth(π)}}
            大阪大学 大学院 工学研究科 平成20年度 院試 数学 問題4 の解答を参照。
    \chapter{積分公式集}
        \section{\texorpdfstring{$\integrate{a}{b}{2x/\sqrt{(x^2-a^2)(b^2-x^2)}}{}{x}=\pi\;(0<a<b)$}{∫\_a\string^b 2x/sqrt((x\string^2-a\string^2)(b\string^2-x\string^2)) = π (0<a<b)}}
            \begin{proof}
                \quad\par
                求めたい積分値を $S$ とする。
                変数変換 $x=a+t$ により次式を得る。
                \[ S = \integrate{0}{b-a}{\frac{2(a+t)}{\sqrt{(t^2+2at)\parens*{(b^2-a^2)-(t^2+2at)}}}}{}{t} \]
                変数変換 $t^2+2at=u$ により次式を得る。
                \[ S = \integrate{0}{b^2-a^2}{\frac{1}{\sqrt{u\parens*{(b^2-a^2)-u}}}}{}{u} \]
                変数変換 $u=(b^2-a^2)v$ により次式を得る。
                \[ S = \integrate{0}{1}{\frac{1}{\sqrt{v(1-v)}}}{}{v} = \Beta{1/2}{1/2} = \pi \]
                （$v=\sin^2\theta$ とおくことで $\Beta{1/2}{1/2} = \pi$ が解る。）
            \end{proof}